
\subsubsection{4.1 Dataset}

All experiments use \textbf{ModelZooDataset CIFAR10-GS} (Unterthiner et al., 2020; Schürholt et al., 2022), available at Zenodo (ID: 6620868). The dataset contains 100 CNN models trained on CIFAR-10 under varied hyperparameters (learning rate, batch size, weight decay, optimizer). Each model is a 3-conv CNN with the following architecture (from h-e1 validation):

\begin{itemize}
\item Input: $28\times 28\times 3$ (grayscale-to-RGB)
\item Conv($3\rightarrow 8, 5\times 5$, no padding) $\rightarrow$ MaxPool(2) $\rightarrow$ LeakyReLU
\item Conv($8\rightarrow 6, 5\times 5$, no padding) $\rightarrow$ MaxPool(2) $\rightarrow$ LeakyReLU
\item Conv($6\rightarrow 4, 2\times 2$, no padding) $\rightarrow$ LeakyReLU $\rightarrow$ Flatten $\rightarrow$ 36 features
\item FC($36\rightarrow 20) \rightarrow$ LeakyReLU $\rightarrow$ FC($20\rightarrow 10)$
\end{itemize}

\begin{table}[htbp]
\centering
\resizebox{\columnwidth}{!}{%
\begin{tabular}{ll}
\toprule
Property & Value \\
\midrule
Models & 100 CNNs \\
Architecture & 3-conv (C: $3\rightarrow 8\rightarrow 6\rightarrow 4)$, 2 FC layers \\
Permutation group & $S_8 \times S_6 \times S_4$ \\
Label & Test accuracy on CIFAR-10 \\
\bottomrule
\end{tabular}
}
\end{table}

This benchmark enables principled comparison with Unterthiner et al. (2020) and Zhou et al. (2023).

\subsubsection{4.2 Encoders}

\begin{table}[htbp]
\centering
\resizebox{\columnwidth}{!}{%
\begin{tabular}{llll}
\toprule
ID & Name & Invariance & Reference \\
\midrule
C0 & Per-layer statistics ($\hat{W}_L$) & Approximately invariant & Unterthiner et al. (2020) \\
C1 & CISE (sinusoidal PE) & Non-invariant & Representative non-invariant baseline (this work) \\
C2 & DeepSets doubly-invariant sum pooling & Exactly invariant & Zaheer et al. (2017) \\
C3 & NFN equivariant & Near-invariant & Zhou et al. (2023) \\
\bottomrule
\end{tabular}
}
\end{table}

C0 and C1 are baselines; C2 is the primary treatment; C3 provides a second architecturally invariant encoder. See Section 4.4 for an implementation note regarding C3.

\subsubsection{4.3 Evaluation Metrics}

\textbf{Primary:} OrbitVar (mean over N = 100 models, K = 50 permutations, float64 precision); MSE\_perm/MSE\_total (ratio for CISE); $R^2$ and Kendall's $\tau$ on held-out testset.

\textbf{Diagnostic:} $R^{2}(C1_avg)$ --- $R^2$ using orbit-averaged CISE embeddings; closure = |$\Delta MSE(C1\rightarrow C2) - MSE_perm^{C1}$| / $MSE_perm^{C1}$ (measures additive decomposition accuracy).

\textbf{Statistical significance:} Wilcoxon signed-rank test on per-model OrbitVar (100 paired observations, C1 vs C2 and C1 vs C3). No multiple-testing correction is applied, as each test addresses a distinct hypothesis.

\subsubsection{4.4 Implementation Note: NFN Library}

The NFN library (\texttt{pip install git+https://github.com/AllanYangZhou/nfn.git}) failed to install in the experimental environment due to dependency conflicts. In h-m3, C3 used DeepSets (C2) as a fallback encoder; consequently, C3 results in Section 5.3 are identical to C2. NFN OrbitVar = 8.905e-08 (reported in h-e1, Table 1) is valid --- the OrbitVar measurement used a simpler NFN instantiation that was successfully installed. NFN downstream $R^2$ comparison against DeepSets remains future work.

\subsubsection{4.5 Compute}

All experiments run on CPU. OrbitVar measurement requires 100 $\times$ 50 = 5,000 forward passes per encoder. LightGBM training takes under 2 minutes per encoder. No GPU was required.

